% Part: first-order-logic
% Chapter: natural-deduction
% Section: proof-theoretic-notions

\documentclass[../../../include/open-logic-section]{subfiles}

\begin{document}

\iftag{FOL}
      {\olfileid{fol}{ntd}{ptn}}
      {\olfileid{pl}{ntd}{ptn}}

\olsection{প্রমাণতাত্ত্বিক ধারণা}

\begin{editorial}
এই অংশে স্বাভাবিক নিষ্পাদনের প্রমাণযোগ্যতা-সম্পর্ক ও সঙ্গতির সংজ্ঞাগুলি
একত্র করা হয়েছে।
\end{editorial}

\begin{explain}
যেভাবে আমরা কয়েকটি গুরুত্বপূর্ণ অর্থতাত্ত্বিক ধারণা (বৈধতা,
অনুসিদ্ধান্ত, পরিতৃপ্তিযোগ্যতা) সংজ্ঞায়িত করেছি, এবার তেমনি তাদের
অনুরূপ \emph{প্রমাণতাত্ত্বিক ধারণা} সংজ্ঞায়িত করি। এগুলি কোনো
!!{structure}s-তে !!{sentence}s-এর পরিতৃপ্তির সাহায্যে সংজ্ঞায়িত নয়;
বরং অন্যগুলি থেকে কয়েকটি !!{sentence}s-এর !!{derivability} বা
!!{nonderivability}-র সাহায্যে সংজ্ঞায়িত। এই ধারণাগুলি যে পরস্পরের
সঙ্গে মিলে যায়, তা এক গুরুত্বপূর্ণ আবিষ্কার। তাদের মিলই
\emph{বিশুদ্ধতা} ও \emph{পূর্ণতা উপপাদ্য}-র বক্তব্য।
\end{explain}


\begin{defn}[উপপাদ্য]
কোনো !!{sentence}~$!A$ একটি \emph{উপপাদ্য}, যদি স্বাভাবিক
নিষ্পাদনে~$!A$-এর এমন কোনো !!a{derivation} থাকে যেখানে সব অনুমিতি
!!{discharged}। $!A$ উপপাদ্য হলে $\Proves !A$ এবং উপপাদ্য না হলে
$\Proves/ !A$ লিখি।
\end{defn}

\begin{defn}[!!^{derivability}]
কোনো !!^a{sentence}~$!A$ !!{sentence}s-এর সেট~$\Gamma$ \emph{থেকে
!!{derivable}}, অর্থাৎ $\Gamma \Proves !A$, যদি সিদ্ধান্ত~$!A$-যুক্ত
এমন একটি !!{derivation} থাকে যার প্রত্যেক অনুমিতি হয় !!{discharged},
নয়তো~$\Gamma$-তে আছে। $!A$ যদি~$\Gamma$ থেকে !!{derivable} না হয়,
তবে $\Gamma \Proves/ !A$ লিখি।
\end{defn}

\begin{defn}[সঙ্গতি]
!!{sentence}s-এর কোনো সেট~$\Gamma$ ঠিক তখনই \emph{অসঙ্গত}, যখন
$\Gamma \Proves \lfalse$। $\Gamma$ অসঙ্গত না হলে, অর্থাৎ $\Gamma
\Proves/ \lfalse$ হলে, সেটিকে \emph{সঙ্গত} বলি।
\end{defn}

\begin{prop}[প্রতিবিম্ব ধর্ম]
\ollabel{prop:reflexivity}
$!A \in \Gamma$ হলে $\Gamma \Proves !A$।
\end{prop}

\begin{proof}
অনুমিতি~$!A$ নিজেই~$!A$-এর এমন একটি !!a{derivation}, যার প্রত্যেক
!!{undischarged} অনুমিতি (অর্থাৎ~$!A$) $\Gamma$-তে আছে।
\end{proof}


\begin{prop}[একঘেয়েতা]
\ollabel{prop:monotonicity}
$\Gamma \subseteq \Delta$ এবং $\Gamma \Proves !A$ হলে $\Delta
\Proves !A$।
\end{prop}

\begin{proof}
$\Gamma$ থেকে~$!A$-এর যেকোনো !!{derivation} আবার~$\Delta$
থেকে~$!A$-এরও !!a{derivation}।
\end{proof}

\begin{prop}[পরিযায়িতা]
\ollabel{prop:transitivity}
$\Gamma \Proves !A$ এবং $\{!A\} \cup \Delta \Proves
!B$ হলে $\Gamma \cup \Delta \Proves !B$।
\end{prop}

\begin{proof}
$\Gamma \Proves !A$ হলে~$\Gamma$-তে সব !!{undischarged} অনুমিতি-যুক্ত
$!A$-এর কোনো !!a{derivation}~$\delta_0$ আছে। $\{!A\} \cup
\Delta \Proves !B$ হলে~$\{!A\} \cup \Delta$-তে সব !!{undischarged}
অনুমিতি-যুক্ত~$!B$-এর কোনো !!a{derivation}~$\delta_1$ আছে। এবার
বিবেচনা করো:
\begin{prooftree}
  \AxiomC{$\Delta, \Discharge{!A}{1}$}
  \RightLabel{$\delta_1$}
  \DeduceC{$!B$}
  \DischargeRule{\Intro{\lif}}{1}
  \UnaryInfC{$!A \lif !B$}
  \AxiomC{$\Gamma$}
  \RightLabel{$\delta_0$}
  \DeduceC{$!A$}
  \RightLabel{\Elim{\lif}}
  \BinaryInfC{$!B$}
\end{prooftree}
এখন সব !!{undischarged} অনুমিতি $\Gamma \cup \Delta$-তে আছে; তাই
এটি $\Gamma \cup \Delta \Proves !B$ দেখায়।
\end{proof}

$\Gamma = \{!A_1, !A_2, \ldots, !A_k\}$ একটি সসীম সেট হলে $\Gamma
\Proves !B$-এর জন্য সরলীকৃত সংকেত $!A_1,!A_2,\ldots,!A_k \Proves
!B$ ব্যবহার করতে পারি। বিশেষত $!A \Proves !B$ মানে
$\{!A\} \Proves !B$।

লক্ষ করো, $\Gamma \Proves !A$ এবং $!A \Proves !B$ হলে $\Gamma
\Proves !B$। এর থেকে আরও পাওয়া যায়: $!A_1, \dots, !A_n \Proves
!B$ এবং প্রত্যেক~$i$-এর জন্য $\Gamma \Proves !A_i$ হলে $\Gamma
\Proves !B$।

\begin{prop}
\ollabel{prop:incons}
নিচের কথাগুলি সমতুল্য।
    \begin{enumerate}
      \item \( \Gamma \) অসঙ্গত।
      \item প্রত্যেক !!{sentence}~\( {!A} \)-এর জন্য \( \Gamma \Proves {!A} \)।
      \item কোনো !!{sentence}~\( {!A} \)-এর জন্য \( \Gamma \Proves {!A} \) এবং \( \Gamma \Proves \lnot {!A} \)।
    \end{enumerate}
\end{prop}

\begin{proof}
অনুশীলনী।
\end{proof}

\tagprob{FOL}
\begin{prob}
\olref[fol][ntd][ptn]{prop:incons} প্রমাণ করো।
\end{prob}
\tagendprob

\tagprob{notFOL}
\begin{prob}
\olref[pl][ntd][ptn]{prop:incons} প্রমাণ করো।
\end{prob}
\tagendprob

\begin{prop}[সংহতি]
  \ollabel{prop:proves-compact}
  \begin{enumerate}
  \item $\Gamma \Proves !A$ হলে এমন একটি সসীম উপসেট $\Gamma_0
    \subseteq \Gamma$ আছে যেন $\Gamma_0 \Proves !A$।
  \item $\Gamma$-র প্রত্যেক সসীম উপসেট সঙ্গত হলে $\Gamma$ সঙ্গত।
  \end{enumerate}
\end{prop}

\begin{proof}
  \begin{enumerate}
    \item $\Gamma \Proves !A$ হলে~$\Gamma$ থেকে~$!A$-এর কোনো
      !!a{derivation}~$\delta$ আছে। $\delta$-র !!{undischarged}
      অনুমিতিগুলির সেটকে~$\Gamma_0$ ধরি। প্রত্যেক !!{derivation}
      সসীম বলে~$\Gamma_0$-তে শুধু সসীমসংখ্যক !!{sentence}s থাকতে পারে।
      তাই $\delta$ কোনো সসীম~$\Gamma_0 \subseteq \Gamma$ থেকে~$!A$-এর
      !!a{derivation}।
    \item এটি (1)-এর বিশেষ ক্ষেত্র $!A \ident \lfalse$-এর
      বিপরীত-প্রতিজ্ঞা।
  \end{enumerate}
\end{proof}

\end{document}
